% 86-68(86쪽) — 두 점 P, Q의 위치 f(t), g(t) 그래프. 스캔 화질이 낮아 TikZ 로 다시 그림.
% 원문에 있는 것만: 두 곡선과 식 라벨 y=f(t), y=g(t) · a·b·c·d 의 세로 파선 · 라벨 y, O, a, b, c, d, e, t.
% 눈금·수치는 원문에 없으므로 넣지 않는다.
% 2026-10-02 전수 검수: f 의 t축 절편이 b 와 c 사이(2.75)에 있었으나, 원본 확대 크롭에서 절편이 c 의 세로 파선과 정확히 같은 자리다.
% t=c 에서 P가 원점에 있다는 원문의 특징이 사라지는 자리라 f 의 좌표열을 다시 잡았다.
\begin{tikzpicture}[x=0.50cm, y=0.50cm, line width=0.55pt]
  \draw[->, >=stealth, line width=0.7pt] (-0.80,0) -- (7.00,0) node[below=1pt, font=\small] {$t$};
  \draw[->, >=stealth, line width=0.7pt] (0,-2.75) -- (0,3.10) node[left=1pt, font=\small] {$y$};
  \draw[densely dashed, line width=0.4pt] (0.95,0) -- (0.95,2.22) (2.20,0) -- (2.20,1.50)
    (3.30,0) -- (3.30,1.63) (4.80,0) -- (4.80,-2.17);
  % y=f(t)
  \draw[line width=0.6pt] plot[smooth, tension=0.7] coordinates {
    (0,0) (0.40,1.25) (0.72,1.95) (1.10,2.32) (1.55,2.40) (1.90,2.12) (2.20,1.50)
    (2.70,0.80) (3.30,0) (3.80,-1.05) (4.30,-1.85) (4.80,-2.17)
    (5.20,-1.85) (5.55,-1.05) (5.90,0) (6.20,0.95) (6.50,2.05) };
  % y=g(t)
  \draw[line width=0.6pt] plot[smooth, tension=0.7] coordinates {
    (0,0) (0.40,0.58) (0.95,1.05) (1.50,1.32) (2.20,1.50) (2.80,1.60) (3.30,1.63)
    (3.90,1.50) (4.50,1.15) (5.10,0.70) (5.55,0.32) (5.90,0) (6.30,-0.52) (6.70,-1.05) };
  \node[below left=-2pt and -3pt, font=\small] at (0,0) {$\pt{O}$};
  \node[below=1pt, font=\small] at (0.95,0) {$a$};
  \node[below=1pt, font=\small] at (2.20,0) {$b$};
  \node[above=1pt, font=\small] at (3.30,0) {$c$};
  \node[below=1pt, font=\small] at (4.80,0) {$d$};
  \node[below=1pt, font=\small] at (5.90,0) {$e$};
  \node[anchor=west, font=\small] at (4.45,2.72) {$y=f(t)$};
  \node[anchor=west, font=\small] at (5.35,-2.15) {$y=g(t)$};
\end{tikzpicture}
